\textbf{Phần 1. Trắc nghiệm lựa chọn}
\setcounterpageref{ex}{0}

\begin{ex} %[5D7V37]
Trên bản đồ tỉ lệ $1 : 50\,000$, quãng đường từ thị trấn A đến thị trấn B đo được là $8\text{ cm}$. Quãng đường thực tế ngoài đời dài bao nhiêu ki-lô-mét?
\choice
{$0{,}4\text{ km}$}
{\True $4\text{ km}$}
{$40\text{ km}$}
{$400\text{ km}$}
\loigiai{
Quãng đường thực tế là:
$8 \times 50\,000 = 400\,000\text{ (cm)}$.\\
Đổi: $400\,000\text{ cm} = 4\text{ km}$.\\
\textbf{Đáp án của câu hỏi là B}
}
\end{ex}

\begin{ex} %[5D7V37]
Một mảnh đất hình chữ nhật có chiều dài thực tế $80\text{ m}$ và chiều rộng thực tế $50\text{ m}$. Trên bản đồ tỉ lệ $1 : 1\,000$, diện tích của mảnh đất đó bằng bao nhiêu xăng-ti-mét vuông?
\choice
{$4\text{ cm}^2$}
{\True $40\text{ cm}^2$}
{$400\text{ cm}^2$}
{$4\,000\text{ cm}^2$}
\loigiai{
Đổi: $80\text{ m} = 8\,000\text{ cm}$; $50\text{ m} = 5\,000\text{ cm}$.\\
Chiều dài trên bản đồ: $8\,000 : 1\,000 = 8\text{ (cm)}$.\\
Chiều rộng trên bản đồ: $5\,000 : 1\,000 = 5\text{ (cm)}$.\\
Diện tích trên bản đồ là: $8 \times 5 = 40\text{ (cm}^2\text{)}$.\\
\textbf{Đáp án của câu hỏi là B}
}
\end{ex}

\begin{ex} %[5D7V37]
Khoảng cách giữa hai địa điểm trên thực tế là $15\text{ km}$. Trên một bản đồ, khoảng cách giữa hai địa điểm đó đo được là $3\text{ cm}$. Tỉ lệ của bản đồ đó là
\choice
{$1 : 50\,000$}
{\True $1 : 500\,000$}
{$1 : 5\,000\,000$}
{$1 : 45\,000$}
\loigiai{
Đổi: $15\text{ km} = 1\,500\,000\text{ cm}$.\\
Tỉ lệ của bản đồ là:
$3 : 1\,500\,000 = 1 : 500\,000$.\\
\textbf{Đáp án của câu hỏi là B}
}
\end{ex}

\begin{ex} %[5D7V37]
Một cái hồ nước hình tròn trên bản đồ tỉ lệ $1 : 2\,000$ có bán kính đo được là $1\text{ cm}$. Bán kính thực tế của hồ nước đó là
\choice
{$2\text{ m}$}
{\True $20\text{ m}$}
{$200\text{ m}$}
{$2\,000\text{ m}$}
\loigiai{
Bán kính thực tế của hồ nước là:
$1 \times 2\,000 = 2\,000\text{ (cm)} = 20\text{ m}$.\\
\textbf{Đáp án của câu hỏi là B}
}
\end{ex}

\begin{ex} %[5D7V37]
Trên bản đồ tỉ lệ $1 : 10\,000$, chu vi của một khu đất hình vuông đo được là $16\text{ cm}$. Diện tích thực tế của khu đất đó bằng bao nhiêu héc-ta?
\choice
{$4\text{ ha}$}
{\True $16\text{ ha}$}
{$64\text{ ha}$}
{$160\text{ ha}$}
\loigiai{
Độ dài cạnh khu đất trên bản đồ là:
$16 : 4 = 4\text{ (cm)}$.\\
Độ dài cạnh thực tế là:
$4 \times 10\,000 = 40\,000\text{ (cm)} = 400\text{ m}$.\\
Diện tích thực tế của khu đất là:
$400 \times 400 = 160\,000\text{ (m}^2\text{)} = 16\text{ ha}$.\\
\textbf{Đáp án của câu hỏi là B}
}
\end{ex}

\begin{ex} %[5D7V37]
Một tuyến đường sắt dài $180\text{ km}$. Khi vẽ trên bản đồ tỉ lệ $1 : 1\,500\,000$, chiều dài tuyến đường sắt trên bản đồ là
\choice
{$1{,}2\text{ cm}$}
{\True $12\text{ cm}$}
{$15\text{ cm}$}
{$120\text{ cm}$}
\loigiai{
Đổi: $180\text{ km} = 18\,000\,000\text{ cm}$.\\
Chiều dài trên bản đồ là:
$18\,000\,000 : 1\,500\,000 = 12\text{ (cm)}$.\\
\textbf{Đáp án của câu hỏi là B}
}
\end{ex}

\begin{ex} %[5D7V37]
Một sân bóng hình chữ nhật trên bản đồ tỉ lệ $1 : 500$ có chiều dài $18\text{ cm}$ và chiều rộng $12\text{ cm}$. Chu vi thực tế của sân bóng đó là
\choice
{$150\text{ m}$}
{\True $300\text{ m}$}
{$600\text{ m}$}
{$3\,000\text{ m}$}
\loigiai{
Chiều dài thực tế là:
$18 \times 500 = 9\,000\text{ (cm)} = 90\text{ m}$.\\
Chiều rộng thực tế là:
$12 \times 500 = 6\,000\text{ (cm)} = 60\text{ m}$.\\
Chu vi thực tế là:
$(90 + 60) \times 2 = 300\text{ (m)}$.\\
\textbf{Đáp án của câu hỏi là B}
}
\end{ex}

\begin{ex} %[5D7V37]
Khi thu nhỏ một hình vuông theo tỉ lệ $1 : 50$, diện tích của hình vuông thu nhỏ so với diện tích hình vuông thực tế giảm đi bao nhiêu lần?
\choice
{$50\text{ lần}$}
{$100\text{ lần}$}
{\True $2\,500\text{ lần}$}
{$5\,000\text{ lần}$}
\loigiai{
Khi kích thước các cạnh giảm đi $50$ lần thì diện tích giảm đi:
$50 \times 50 = 2\,500\text{ (lần)}$.\\
\textbf{Đáp án của câu hỏi là C}
}
\end{ex}

\begin{ex} %[5D7V37]
Bản đồ phòng học vẽ theo tỉ lệ $1 : 100$. Chiều dài phòng học đo được trên bản đồ là $9\text{ cm}$, chiều rộng là $7\text{ cm}$. Diện tích thực tế của phòng học là
\choice
{$16\text{ m}^2$}
{\True $63\text{ m}^2$}
{$630\text{ m}^2$}
{$6\,300\text{ m}^2$}
\loigiai{
Chiều dài thực tế là: $9 \times 100 = 900\text{ cm} = 9\text{ m}$.\\
Chiều rộng thực tế là: $7 \times 100 = 700\text{ cm} = 7\text{ m}$.\\
Diện tích thực tế là: $9 \times 7 = 63\text{ (m}^2\text{)}$.\\
\textbf{Đáp án của câu hỏi là B}
}
\end{ex}

\begin{ex} %[5D7V37]
Trên bản đồ tỉ lệ $1 : 250\,000$, một cây cầu dài $2\text{ cm}$. Độ dài thật của cây cầu đó là
\choice
{$0{,}5\text{ km}$}
{\True $5\text{ km}$}
{$50\text{ km}$}
{$500\text{ m}$}
\loigiai{
Độ dài thật của cây cầu là:
$2 \times 250\,000 = 500\,000\text{ (cm)} = 5\text{ km}$.\\
\textbf{Đáp án của câu hỏi là B}
}
\end{ex}

\begin{ex} %[5D7V37]
Một thửa ruộng hình tam giác trên bản đồ tỉ lệ $1 : 1\,000$ có cạnh đáy $6\text{ cm}$ và chiều cao $4\text{ cm}$. Diện tích thực tế của thửa ruộng đó là
\choice
{$240\text{ m}^2$}
{\True $1\,200\text{ m}^2$}
{$2\,400\text{ m}^2$}
{$12\,000\text{ m}^2$}
\loigiai{
Cạnh đáy thực tế: $6 \times 1\,000 = 6\,000\text{ cm} = 60\text{ m}$.\\
Chiều cao thực tế: $4 \times 1\,000 = 4\,000\text{ cm} = 40\text{ m}$.\\
Diện tích thực tế: $60 \times 40 : 2 = 1\,200\text{ (m}^2\text{)}$.\\
\textbf{Đáp án của câu hỏi là B}
}
\end{ex}

\begin{ex} %[5D7V37]
Khoảng cách từ Hà Nội đến Hải Phòng là $105\text{ km}$. Trên bản đồ tỉ lệ $1 : 1\,000\,000$, khoảng cách này là
\choice
{$1{,}05\text{ cm}$}
{\True $10{,}5\text{ cm}$}
{$105\text{ cm}$}
{$0{,}105\text{ cm}$}
\loigiai{
Đổi: $105\text{ km} = 10\,500\,000\text{ cm}$.\\
Khoảng cách trên bản đồ là:
$10\,500\,000 : 1\,000\,000 = 10{,}5\text{ (cm)}$.\\
\textbf{Đáp án của câu hỏi là B}
}
\end{ex}

\textbf{Phần 2. Trắc nghiệm đúng sai}
\setcounterpageref{ex}{0}

\begin{ex} %[5D7V37]
Trên bản đồ quy hoạch khu đô thị có tỉ lệ $1 : 2\,000$, một công viên có dạng hình chữ nhật với chiều dài $15\text{ cm}$ và chiều rộng $10\text{ cm}$.
\choiceTF
{\True Chiều dài thực tế của công viên là $300\text{ m}$}
{\True Chiều rộng thực tế của công viên là $200\text{ m}$}
{\True Chu vi thực tế của công viên là $1\,000\text{ m}$}
{\True Diện tích thực tế của công viên là $6\text{ ha}$}
\loigiai{
\begin{itemchoice}
\item Đúng. Chiều dài thật: $15 \times 2\,000 = 30\,000\text{ (cm)} = 300\text{ m}$.
\item Đúng. Chiều rộng thật: $10 \times 2\,000 = 20\,000\text{ (cm)} = 200\text{ m}$.
\item Đúng. Chu vi thật: $(300 + 200) \times 2 = 1\,000\text{ (m)}$.
\item Đúng. Diện tích thật: $300 \times 200 = 60\,000\text{ (m}^2\text{)} = 6\text{ ha}$.
\end{itemchoice}
}
\end{ex}

\begin{ex} %[5D7V37]
Một đoạn đường cao tốc có chiều dài thực tế là $40\text{ km}$.
\choiceTF
{\True Trên bản đồ tỉ lệ $1 : 1\,000\,000$, đoạn đường đó dài $4\text{ cm}$}
{\True Trên bản đồ tỉ lệ $1 : 500\,000$, đoạn đường đó dài $8\text{ cm}$}
{Bản đồ có tỉ lệ càng lớn thì hình vẽ đối tượng trên bản đồ càng nhỏ}
{\True Nếu trên một bản đồ khác đoạn đường đó dài $2\text{ cm}$ thì tỉ lệ của bản đồ đó là $1 : 2\,000\,000$}
\loigiai{
\begin{itemchoice}
\item Đúng. $40\text{ km} = 4\,000\,000\text{ cm} \implies 4\,000\,000 : 1\,000\,000 = 4\text{ (cm)}$.
\item Đúng. $4\,000\,000 : 500\,000 = 8\text{ (cm)}$.
\item Sai. Bản đồ tỉ lệ càng lớn thì hình vẽ trên bản đồ càng chi tiết và to hơn.
\item Đúng. $4\,000\,000 : 2 = 2\,000\,000$ nên tỉ lệ là $1 : 2\,000\,000$.
\end{itemchoice}
}
\end{ex}

\textbf{Phần 3. Trắc nghiệm trả lời ngắn}
\setcounterpageref{ex}{0}

\begin{ex} %[5D7V37]
Trên bản đồ tỉ lệ $1 : 50\,000$, khoảng cách giữa hai điểm A và B là $6\text{ cm}$. Khoảng cách thực tế giữa hai điểm đó là bao nhiêu ki-lô-mét?
\shortans{3}
\loigiai{
Khoảng cách thực tế là:
$6 \times 50\,000 = 300\,000\text{ (cm)}$.\\
Đổi: $300\,000\text{ cm} = 3\text{ km}$.
}
\end{ex}

\begin{ex} %[5D7V37]
Một mảnh đất hình vuông có độ dài cạnh thực tế là $50\text{ m}$. Trên bản đồ tỉ lệ $1 : 500$, chu vi của hình vuông đó là bao nhiêu xăng-ti-mét?
\shortans{40}
\loigiai{
Đổi: $50\text{ m} = 5\,000\text{ cm}$.\\
Cạnh hình vuông trên bản đồ là:
$5\,000 : 500 = 10\text{ (cm)}$.\\
Chu vi hình vuông trên bản đồ là:
$10 \times 4 = 40\text{ (cm)}$.
}
\end{ex}

\begin{ex} %[5D7V37]
Khoảng cách giữa hai thành phố là $160\text{ km}$. Trên bản đồ tỉ lệ $1 : 2\,000\,000$, khoảng cách giữa hai thành phố đó là bao nhiêu xăng-ti-mét?
\shortans{8}
\loigiai{
Đổi: $160\text{ km} = 16\,000\,000\text{ cm}$.\\
Khoảng cách trên bản đồ là:
$16\,000\,000 : 2\,000\,000 = 8\text{ (cm)}$.
}
\end{ex}

\begin{ex} %[5D7V37]
Một hồ nước hình chữ nhật trên bản đồ tỉ lệ $1 : 1\,000$ có chiều dài $5\text{ cm}$ và chiều rộng $3\text{ cm}$. Diện tích thực tế của hồ nước đó là bao nhiêu mét vuông?
\shortans{1500}
\loigiai{
Chiều dài thật: $5 \times 1\,000 = 50\text{ m}$.\\
Chiều rộng thật: $3 \times 1\,000 = 30\text{ m}$.\\
Diện tích thực tế là:
$50 \times 30 = 1\,500\text{ (m}^2\text{)}$.
}
\end{ex}

\begin{ex} %[5D7V37]
Trên bản đồ tỉ lệ $1 : 10\,000$, một con kênh dài $25\text{ cm}$. Chiều dài thật của con kênh đó là bao nhiêu mét?
\shortans{2500}
\loigiai{
Chiều dài thật của con kênh là:
$25 \times 10\,000 = 250\,000\text{ (cm)}$.\\
Đổi: $250\,000\text{ cm} = 2\,500\text{ m}$.
}
\end{ex}

\begin{ex} %[5D7V37]
Khoảng cách giữa hai ngọn hải đăng là $45\text{ km}$. Trên hải đồ, khoảng cách đo được là $9\text{ cm}$. Mẫu số của tỉ lệ xích trên hải đồ có dạng $a00\,000$. Tìm chữ số $a$.
\shortans{5}
\loigiai{
Đổi: $45\text{ km} = 4\,500\,000\text{ cm}$.\\
Mẫu số của tỉ lệ là:
$4\,500\,000 : 9 = 500\,000$.\\
Vậy chữ số $a = 5$.
}
\end{ex}
